
\documentclass[a4paper,10pt]{report}

%\usepackage[landscape]{geometry}
\usepackage[cm-default]{fontspec}
\setromanfont{FreeSerif}
\setsansfont{FreeSans}
\setmonofont{FreeMono}
\usepackage[utf8x]{inputenc}
\usepackage{fontspec}
\usepackage{xunicode}
\usepackage{xltxtra}
\usepackage{xgreek}
\usepackage{amsmath}
\usepackage{unicode-math}
\usepackage{ulem}
\usepackage{color}
\usepackage{verbatim}
\usepackage{nopageno}
\usepackage{graphicx}
\usepackage{textpos}
\setlength{\TPHorizModule}{1cm}
\setlength{\TPVertModule}{1cm}
%\usepackage[colorgrid,texcoord]{eso-pic}
\usepackage[outline]{contour}
\usepackage{wrapfig}
\usepackage{url}
\usepackage{color}
\usepackage{tikz}
\usepackage{fancybox,fancyhdr}
\usepackage{subfigure}
\usepackage{pstricks}
\usepackage{epsfig}
\usepackage{multicol}
\usepackage{listings}
\usepackage{enumerate}
\usepackage{hyperref}
\hypersetup{
  bookmarks=true,
  bookmarksopen=true,
  pdfborder=false,
  pdfpagemode=UseNone,
  raiselinks=true,
  pdfhighlight={/P},
  colorlinks,
  citecolor=black,
  filecolor=black,
  linkcolor=black,
  urlcolor=black
}


\usepackage{fontspec}
\usepackage{xunicode}
\usepackage{xltxtra}


% Margins
%\setlength{\textwidth}{24cm}

%\setlength{\voffset}{0in}
%\setlength{\textheight}{6.8in}
% Colors
\definecolor{cbrown}{rgb}{0.49,0.24,0.07}
\definecolor{cmpez}{rgb}{0.92,0.88,0.79}
\definecolor{cmpez2}{rgb}{0.62,0.33,0.34}
\definecolor{cwhite}{rgb}{1,1,1}
\definecolor{cblack}{rgb}{0,0,0}
\definecolor{cred}{rgb}{0.9,0.15,0.15}
\definecolor{clblue}{rgb}{0.57,0.85,0.97}
\definecolor{corange}{rgb}{0.86,0.49,0.18}
\definecolor{clorange}{rgb}{0.95,0.84,0.61}
\definecolor{cyellow}{rgb}{1,0.95,0.16}
\definecolor{cgreen}{rgb}{0.69,0.8,0.31}
\definecolor{clbrown}{rgb}{0.78,0.67,0.43}
\definecolor{cmagenta}{rgb}{0.79,0.05,0.54}
\definecolor{cgray}{rgb}{0.76,0.73,0.63}
\definecolor{bluesite}{rgb}{0.05,0.43,0.67}
\definecolor{redsite}{rgb}{0.90,0.15,0.15}
\definecolor{backsite}{rgb}{0.07,0.07,0.07}
% ----------- from Costas -----------------------------

%---------------   Main Fonts    -----------------%   

% \setmainfont[Mapping=tex-text]{Calibri}
% \setmainfont[Mapping=tex-text]{Times New Roman}
%\setmainfont[Mapping=tex-text]{Droid Serif}
%\setmainfont[Mapping=tex-text]{Cambria}
%\setmainfont[Mapping=tex-text]{cm-unicode}
% \setmainfont[Mapping=tex-text]{Gentium}
% \setmainfont[Mapping=tex-text]{GFS Didot}
% \setmainfont[Mapping=tex-text]{Comic Sans MS}
 \setmainfont[Mapping=tex-text]{Ubuntu}
% \setmainfont[Mapping=tex-text]{Myriad Pro}
% \setmainfont[Mapping=tex-text]{CMU Concrete}
%\setmainfont[Mapping=tex-text]{DejaVu Sans}
%\setmainfont[Mapping=tex-text]{KerkisSans}
%\setmainfont[Mapping=tex-text]{KerkisCaligraphic}
% \setmainfont[Mapping=tex-text]{Segoe Print}
% \setmainfont[Mapping=tex-text]{Gabriola}

%------------------------------------------------%

% ------------ Mathematics Fonts ----------------%
\setmathfont{Asana-Math.ttf}
% -----------------------------------------------%
% Misc
\author{Κασωτάκη Ε. - Σμαραγδάκης Κ.}
\title{www.math24.gr}
% ---------------- Formation --------------------%


\setlength\topmargin{-0.7cm}
\addtolength\topmargin{-\headheight}
\addtolength\topmargin{-\headsep}
\setlength\textheight{26cm}
\setlength\oddsidemargin{-0.54cm}
\setlength\evensidemargin{-0.54cm}
%\setlength\marginparwidth{1.5in}
\setlength\textwidth{17cm}

\RequirePackage[avantgarde]{quotchap}
\renewcommand\chapterheadstartvskip{\vspace*{0\baselineskip}}
\RequirePackage[calcwidth]{titlesec}
\titleformat{\section}[hang]{\bfseries}
{\Large\thesection}{12pt}{\Large}[{\titlerule[0.9pt]}]
%--------------------------------------------------%
\usepackage{draftwatermark}
\SetWatermarkText{www.math24.gr}
\SetWatermarkLightness{0.9}
\SetWatermarkScale{0.6}
%--------------------------------------------------%


\usepackage{xcolor}
\usepackage{amsthm}
\usepackage{framed}
\usepackage{parskip}

\colorlet{shadecolor}{bluesite!20}

\newtheorem*{orismos}{Ορισμός}

\newenvironment{notation}
  {\begin{shaded}\begin{theorem}}
  {\end{theorem}\end{shaded}}




% Document begins
\begin{document}
\pagestyle{fancy}
\fancyhead{}
\fancyfoot{}
\renewcommand{\headrulewidth}{0pt}
\renewcommand{\footrulewidth}{0pt}


\fancyhead[LO,LE]{
 %\textblockcolor{backsite}
 %\begin{textblock}{5}(-2,-0.55)
  %\rule{0cm}{1cm}
 %\end{textblock}
 \textblockcolor{bluesite}
 \begin{textblock}{5}(-1.5,-0.55)
  \rule{0cm}{1cm}
 \end{textblock}
 %\textblockcolor{bluesite}
 %\begin{textblock}{14}(5.5,-0.55)
 % \rule{0cm}{1cm}
 %\end{textblock}
 \begin{textblock}{0}(-1,-0.25)
 \color{cwhite} \begin{Large}www.math24.gr\end{Large}
 \end{textblock}
\begin{textblock}{0}(16,-0.4)
 \color{cwhite} \includegraphics[height=1.5cm]{math24_logo.png}
 \end{textblock}
\textblockcolor{white}
\begin{textblock}{17}(-1,28)
 \color{backsite} \begin{small}Copyright \textcopyright 2011, Κασωτάκη Ε.(ikasotaki@gmail.com) - Σμαραγδάκης Κ.(kesmarag@gmail.com)\end{small}
 \end{textblock}
}
 
%\begin{shaded}
%\begin{center}
%\huge \textbf{Φυλλάδιο Ασκήσεων}\\
%Πρόσθεση ρητών αριθμών
%\end{center} 
%\textbf{Μαθηματικά Α' Γυμνασίου} \hfill \textbf{Ημερομηνία Παράδοσης : \hspace{2em} }
%\subsection*{Ονοματεπώνυμο :\hfill  \hspace{5em}}
%\end{shaded}
%\vspace{2em}


\section*{Άσκηση 1  (Η έννοια του κλάσματος)\hfill \small{}}
Γνωρίζουμε ότι μια ημέρα έχει 24 ώρες. Να βρείτε ποιο μέρος της ημέρας είναι 
\begin{enumerate}[i)]
 \item οι 2 ώρες
 \item οι 4 ώρες
 \item οι 8 ώρες 
 \item οι 6 ώρες 
 \item οι 12 ώρες
\end{enumerate}
\textbf{Λύση}
\begin{enumerate}[i)]
 \item οι 2 ώρες της ημέρας είναι τα $\dfrac{2}{24}$ της ημέρας
 \item οι 4 ώρες της ημέρας είναι τα $\dfrac{4}{24}$ της ημέρας
 \item οι 8 ώρες της ημέρας είναι τα $\dfrac{8}{24}$ της ημέρας
 \item οι 6 ώρες της ημέρας είναι τα $\dfrac{6}{24}$ της ημέρας
 \item οι 12 ώρες της ημέρας είναι τα $\dfrac{12}{24}$ της ημέρας
\end{enumerate}







\section*{Άσκηση 2 (Σύγκριση κλασμάτων με τη μονάδα) \hfill \small{}}
Να βρείτε ποια από τα παρακάτω κλάσματα είναι μικρότερα, ποια μεγαλύτερα της μονάδα; και ποια είναι ίσα με τη μονάδα;
\begin{multicols}{2}
\begin{enumerate}[1)]
%\begin{itemize}
 \item $\dfrac{1}{5}$
 \item $\dfrac{2}{8}$
 \item $\dfrac{8}{8}$
 \item $\dfrac{7}{3}$
 \item $\dfrac{7}{7}$
 \item $\dfrac{3}{6}$
 \item $\dfrac{2}{3}$
 \item $\dfrac{3}{3}$
 \item $\dfrac{5}{16}$
 \item $\dfrac{7}{2}$
%\end{itemize}
\end{enumerate}
\end{multicols}
\textbf{Λύση}
\begin{multicols}{2}
\begin{enumerate}[1)]
 \item $\dfrac{1}{5}<1$ διότι $1<5$ 
 \item $\dfrac{2}{8}<1$ διότι $2<8$
 \item $\dfrac{8}{8}=1$ διότι $8=8$
 \item $\dfrac{7}{3}>1$ διότι $7>3$
 \item $\dfrac{7}{7}=1$ διότι $7=7$
 \item $\dfrac{3}{6}<1$ διότι $3<6$
 \item $\dfrac{2}{3}<1$ διότι $2<3$
 \item $\dfrac{3}{3}=1$ διότι $3=3$
 \item $\dfrac{5}{16}<1$ διότι $5<16$
 \item $\dfrac{7}{2}>1$ διότι $7>2$
\end{enumerate}
\end{multicols}



\section*{Άσκηση 3 (Υπολογισμός του μέρους από το όλο)\hfill \small{}}
Ένα πακέτο χωράει $20$ τσίχλες. Κάποια στιγμή βρήκαμε ότι το πακέτο ήταν γεμάτο κατά τα $\dfrac{3}{4}$. 
Πόσες τσίχλες περιείχε τότε το πακέτο;\\
\textbf{Λύση}\\ \\
Τα $\dfrac{4}{4}$ του πακέτου χωράνε $20$ τσίχλες.\\
Το $\dfrac{1}{4}$ του πακέτου χωράει $20:4=5$ τσίχλες.\\
Άρα τα $\dfrac{3}{4}$ του πακέτου χωράνε $3\cdot 5=15$ τσίχλες.


\section*{Άσκηση 4 (Υπολογισμός του μέρους από το όλο)\hfill \small{}}
Μία παπουτσοθήκη χωράει $30$ ζευγάρια παπούτσια. Πόσα παπούτσια χωράνε τα  $\dfrac{5}{6}$; \\ \\
\textbf{Λύση}\\
Τα $\dfrac{6}{6}$ της παπουτσοθήκης χωράνε $30$ ζευγάρια παπούτσια.\\
Το $\dfrac{1}{6}$ της παπουτσοθήκης χωράει $30:6=5$ ζευγάρια παπούτσια.\\
Άρα τα $\dfrac{5}{6}$ της παπουτσοθήκης χωράνε $5\cdot 5=25$ ζευγάρια παπούτσια.




\section*{Άσκηση 5 (Υπολογισμός του όλου από το ένα μέρος)\hfill \small{}}
Τα $\dfrac{2}{6}$ του κιλού μοσχάρι κοστίζουν $4\texteuro.$ Πόσο κοστίζουν τα $\dfrac{2}{3}$ του κιλού;\\ \\
\textbf{Λύση}\\
Τα $\dfrac{2}{6}$ του κιλού κοστίζουν $4\texteuro$.\\
Άρα το $\dfrac{1}{6}$ κοστίζει $4\texteuro:2=2\texteuro$ .\\
     Άρα  το ένα κιλό, δηλαδή τα $\dfrac{6}{6}$ κοστίζει $6\cdot 2\texteuro=12\texteuro$.\\
       Το ένα κιλό, δηλαδή τα $\dfrac{3}{3}$ κοστίζει $12\texteuro$.\\
      Άρα το $\dfrac{1}{3}$ του κιλού κοστίζει $12\texteuro :3=4\texteuro$.\\
       Άρα τα $\dfrac{2}{3}$ του κιλού κοστίζουν $2\cdot 4\texteuro=8\texteuro$

\section*{Άσκηση 6 (Υπολογισμός του όλου από το ένα μέρος)\hfill \small{}}
Τα $\dfrac{4}{5}$ ενός ντεπόζιτου χωράνε  $800lt$. Πόσα $lt$ χωράνε τα $\dfrac{9}{10}$;\\ \\
\textbf{Λύση}\\
Τα $\dfrac{4}{5}$ του ντεπόζιτου χωράνε $800lt$.\\
Άρα το $\dfrac{1}{5}$ χωράει $800lt:4=200lt$ .\\
     Άρα  το ένα ντεπόζιτο, δηλαδή τα $\dfrac{5}{5}$ χωράνε $5\cdot 200lt=1000lt$.\\
       Το ένα ντεπόζιτο, δηλαδή τα $\dfrac{10}{10}$ χωράνε $1000lt$.\\
      Άρα το $\dfrac{1}{10}$ του ντεπόζιτου χωράει $1000lt :10=100lt$.\\
       Άρα τα $\dfrac{9}{10}$ του ντεπόζιτου χωράνε $9\cdot 100lt=900lt$
\end{document}
